% =============================================================================
% Section 4 -- The APIP Framework
% Contains Figure 1 (fig:lifecycle), drawn with a plain tabular so that no
% additional LaTeX packages (tikz, etc.) are required.
% =============================================================================
\section{The APIP Framework}
\label{sec:framework}

We now describe the APIP as a structured lifecycle protocol. The framework
consists of five phases: Trigger, Cause Attribution, Intervention Planning,
Remediation Window, and Exit Evaluation. Fig.~\ref{fig:lifecycle} illustrates
the lifecycle; each phase is described in detail below.

\begin{figure}[t]
  \centering
  \footnotesize
  \setlength{\tabcolsep}{3pt}
  \renewcommand{\arraystretch}{1.3}
  \begin{tabular}{@{}|c|c|c|@{}}
    \hline
    \multicolumn{3}{|c|}{1. Trigger} \\
    \hline
    \multicolumn{3}{|c|}{$\downarrow$} \\
    \multicolumn{3}{|c|}{2. Cause Attribution} \\
    \hline
    \multicolumn{3}{|c|}{$\downarrow$} \\
    \multicolumn{3}{|c|}{3. Intervention Plan} \\
    \hline
    \multicolumn{3}{|c|}{$\downarrow$} \\
    \multicolumn{3}{|c|}{4. Remediation Window} \\
    \hline
    \multicolumn{3}{|c|}{$\downarrow$} \\
    \multicolumn{3}{|c|}{5. Exit Evaluation} \\
    \hline
  \end{tabular}
  \caption{The five-phase APIP lifecycle.}
  \label{fig:lifecycle}
\end{figure}

\subsection{Phase 1: Trigger}
\label{subsec:trigger}

An APIP is initiated when an agent's performance crosses a pre-specified
trigger threshold on one or more monitored metrics. This presupposes the
existence of a \emph{behavioral contract}---a formal specification of the
performance envelope the agent is expected to maintain at deployment time. The
behavioral contract is a deliberate design artifact that most teams currently
skip, delegating performance expectations to informal organizational memory.

A minimal behavioral contract specifies task surface, success metrics and
measurement methodology, acceptable performance ranges, escalation policy, and
human review integration points. Trigger conditions are threshold crossings on
these metrics---e.g., task completion rate below 82\% over a 7-day window.

Two trigger modalities are relevant: absolute threshold triggers, which fire
when a metric crosses a fixed floor, and relative drift triggers, which fire
when a metric declines significantly from its recent baseline regardless of
absolute level. Both are necessary; the former catches severe degradation, the
latter catches gradual drift that may not breach an absolute floor for some
time.

\subsection{Phase 2: Cause Attribution}
\label{subsec:cause-attribution}

Once triggered, the APIP requires a structured cause attribution process before
any intervention is applied. The goal is to identify which failure mode
category (Section~\ref{sec:taxonomy}) best explains the degradation. Tier
mismatch from misdiagnosis is the primary source of avoidable regression: prompt
engineering applied to a retrieval freshness problem, for instance, consumes
engineering cycles with no improvement.

This attribution problem is distinct from, and complementary to, recent work on
automated failure attribution in agentic systems. That literature localizes the
responsible agent or step within a single failed trajectory. Causal-inference
frameworks combine Shapley-based agent-level blame assignment with causal
discovery over execution logs to identify critical failure steps, motivated by
the finding that state-of-the-art methods achieve under 15\% accuracy at
locating root-cause steps on the Who\&When
benchmark~\cite{ma2025attribution,zhang2025whichagent}. Lightweight causal-graph
tracing supports post-hoc root-cause localization in deployed multi-agent
workflows~\cite{agenttrace2026}. APIP Phase 2 operates one level of aggregation
above: it classifies a window-level degradation into a remediation-oriented
failure-mode category to select an intervention tier. The two levels compose
rather than compete. Trajectory-level attribution outputs are natural evidence
inputs to Phase 2, since recurring step- or agent-level attributions across
sampled failures constrain the failure locus and the plausible failure-mode
classification. Recent root-cause taxonomies make the point explicitly:
localizing the failing step falls short of diagnosing the underlying cause,
whether poor prompt design, improper task decomposition, or logical
deadlock~\cite{ma2025lifecycle}. That diagnostic gap is what the APIP's
mandatory taxonomy classification is designed to close at the deployment level.

The cause attribution process should include: (1) disaggregated performance
analysis by input category, user segment, and topic cluster to isolate failure
locus; (2) analysis of a sampled set of failure cases, applying automated
trajectory-level attribution methods where tooling
permits~\cite{ma2025attribution,agenttrace2026,zhu2026agentdebugx}, with a
qualified human reviewer adjudicating the resulting evidence; (3) retrieval
audit if behavioral drift is suspected; (4) tool invocation log analysis if tool
misuse is suspected; and (5) a structured classification of the failure into the
four-category taxonomy. The phase closes with a written cause attribution report
recording the evidence examined, any automated attribution outputs, and the
adjudicated classification; this report becomes part of the APIP record.

\subsection{Phase 3: Intervention Planning and the Tier Model}
\label{subsec:tiers}

Interventions are organized into three tiers of increasing cost, disruption,
and required lead time. The cause attribution output selects the appropriate
starting tier; escalation to higher tiers is permitted if lower tiers prove
insufficient within the remediation window.

\textbf{Tier 1 --- Prompt-Level Interventions.} Prompt-level interventions
modify the system prompt, few-shot examples, or instruction set: constraint
injection, persona adjustment, example replacement, context restructuring.
Tier 1 is fastest to implement, least disruptive, and easiest to roll back. It
is the appropriate first response to alignment creep and many instances of
input brittleness.

\textbf{Tier 2 --- Retrieval and Context Interventions.} Tier 2 interventions
modify the agent's information environment: updating RAG indices, revising
knowledge bases, modifying tool schemas, adjusting context injection pipelines.
They address behavioral drift caused by world-state divergence and tool misuse
from stale tool descriptions. They carry moderate cost and require contract
validation before re-deployment.

\textbf{Tier 3 --- Model-Level Interventions.} Tier 3 interventions modify the
underlying model: fine-tuning, DPO, or replacement. These carry the highest
cost and regression risk. They are warranted when lower tiers are exhausted,
when the failure mode is embedded in model weights, or when the overall
performance profile cannot be returned to contract range through prompting or
retrieval adjustments.

\subsection{Phase 4: Remediation Window and Check-In Cadence}
\label{subsec:window}

The remediation window is a bounded time period during which the APIP
interventions are implemented and their effects are evaluated. A defined window
serves two functions: it creates urgency for the remediation process, and it
provides a structured decision point at which the APIP outcome is formally
evaluated. The duration of the window should be calibrated to the intervention
tier---Tier 1 interventions may warrant a 7--14 day window, while Tier 3
interventions may require 30--60 days to implement, validate, and stabilize.
This mirrors the organizational PIP practice, where durations typically range
from 30 to 90 days, calibrated to the complexity of the performance
gap~\cite{kirkpatrick2016four}.

Check-ins are scheduled evaluation points within the window at which the
agent's performance on APIP-relevant metrics is reviewed against the recovery
trajectory. Following Kirkpatrick's~\cite{kirkpatrick1994evaluating}
hierarchical evaluation principle, each check-in should assess improvement at
multiple levels: immediate output quality (Level 1/2), behavioral pattern
change on the failure-mode class (Level 3), and business outcome metrics
(Level 4). An intervention that improves surface metrics without producing
Level 3/4 improvement should be flagged as potentially superficial. Check-ins
are early warning mechanisms: no improvement at the first check-in triggers
tier escalation rather than waiting for window expiry. A recommended cadence is
three check-ins at 25\%, 50\%, and 75\% of elapsed time.

\subsection{Phase 5: Exit Evaluation}
\label{subsec:exit}

At the close of the remediation window, a formal exit evaluation determines one
of three outcomes: \emph{Resolution} (the agent has returned to within the
behavioral contract envelope and the APIP is closed), \emph{Extension} (the
agent shows improvement but has not yet reached the threshold; the APIP window
is extended with modified parameters), or \emph{Escalation} (the agent has not
recovered and more fundamental intervention---including potential model
replacement or agent retirement---is warranted). The evaluation record becomes
part of the organizational AI audit trail---providing governance documentation
and institutional memory for future deployments.
